\documentclass[dvipdfmx]{jarticle}
\pagestyle{empty}
\usepackage{tikz}
\usetikzlibrary{shapes,calc}
\newcommand{\syntacticObject}[3]{\node[draw,rectangle,rounded corners=5pt,thick,inner sep=5pt] (#2) at #3 {{\hbox to 10pt{}\bf {#1}\hbox to 10pt{}}}}
\newcommand{\processObject}[3]{\node[draw,rectangle,thick,inner sep=5pt] (#2) at #3 {\hbox to 10pt{}{\bf #1}\hbox to 10pt{}}}
\begin{document}
\begin{tikzpicture}
  \syntacticObject{ソースプログラム}{SRC}{(0,0)};
  \processObject{字句解析系}{LEX}{(0,-1)};
  \syntacticObject{トークン列}{TOKEN}{(0,-2)};
  \processObject{構文解析系}{PAR}{(0,-3)};
  \syntacticObject{解析木}{TREE}{(0,-4)};
  \processObject{中間コード生成系}{ILGEN}{(0,-5)};
  \syntacticObject{中間コード列}{ILSEQ}{(0,-6)};
  \processObject{コード最適化}{OPT}{(0,-7)};
  \processObject{コード生成系}{CODEGEN}{(0,-8)};
  \syntacticObject{目的コード}{OBJCODE}{(0,-9)};
  \path[->,thick] (SRC) edge (LEX);
  \path[->,thick] (LEX) edge (TOKEN);
  \path[->,thick] (TOKEN) edge (PAR);
  \path[->,thick] (PAR) edge (TREE);
  \path[->,thick] (TREE) edge (ILGEN);
  \path[->,thick] (ILGEN) edge (ILSEQ);
  \path[->,thick] (ILSEQ) edge (OPT);
  \path[->,thick] (OPT) edge (CODEGEN);
  \path[->,thick] (CODEGEN) edge (OBJCODE);
\end{tikzpicture}
\end{document}
